\section{Slices keep logical coordinates}
\label{sec:sequence-sql-slicing}

Postgres support array subscription operations natively, so I decided to support slicing as well in first version of BioseqDB.

A slice selects positions \([a,b)\): it includes \(a\) and excludes \(b\).
BioSeqDB uses zero-based positions, so \(S[0]\) is the first symbol of the slice.

\warn{0-based indexing}{
  This contradicts the fact that Postgres slices are 1-based, but it wasn't intentional decision. I had limited experience with Postgres array functionality and this realisation came only after most of the support was already implemented.
  What is more interesting that SeqKit - tool for operating on FASTA/Q files use 1-based indexing~\cite{seqkit_subseq_usage}.
  This incompatibility might be something to reconsider in the future.
}

I provided utility functions for common operations like reversing and reverse-complementing slices as well, example available in listing~\ref{lst:sequence-sql-examples}.
\begin{lstlisting}[style=seq-nine-sql,
  caption={Zero-based slices and sequence transforms in SQL.},
  label={lst:sequence-sql-examples}]
WITH v(s) AS (
  VALUES ('ACGTAACCGTTA'::bioseq.seq('dna4', 32))     -- ACGTAACCGTTA
)
SELECT s[2:10]::text AS middle,                       -- GTAACCGT
       ((s[2:10])[1:5])::text AS nested,              -- TAAC
       s[-4:]::text AS suffix,                        -- GTTA
       bioseq.reverse(s[2:10])::text AS rev,          -- TGCCAATG
       bioseq.reverse_complement(s[2:10])::text AS rc -- ACGGTTAC
FROM v;
\end{lstlisting}

The bounds follow a small set of intuitive rules simiar to design of slices in Python~\cite{python_common_sequence_operations}:
\begin{itemize}
  \item A missing lower bound means zero. A missing upper bound means the
        sequence length.
  \item A negative bound counts from the end (example: for length 12, \(-4\) means 8) and bounds outside the sequence are clamped to its ends.
  \item An upper bound below the lower bound gives an empty slice, so \texttt{s[4:2]} has length zero.
  \item User can extract any element (symbol) from slice, which has separate SQL type \texttt{LETTER}. An out-of-range element, such as \texttt{s[12]} in this example,
        gives SQL \texttt{NULL}. An explicit null slice bound also gives \texttt{NULL} i.e $s[a:\text{NULL}]$ is equivalent to $s[\text{NULL}:a]$ and is also equivalent to $s[\text{NULL}:\text{NULL}]$ (all of them result in \texttt{NULL})
\end{itemize}

Slicing can be implemented by just holding slice coordinates, without actually loading any sequences, using very simple coordinate arithmetic shown in the figure below~\ref{fig:sequence-slice-coordinates}.
For a slice view \(V=S[a,b)\), view position \(j\) refers to
\(S[a+j]\). A second slice can combine its offsets directly to get:
\[
  V[c,d)=S[a+c,a+d).
\]
The example's \(S[2,10)[1,5)\) is \(S[3,7)\), which reads
\texttt{TAAC}. The reversal operation is represented by reversing the start and end coordinates. For
\(W=\operatorname{reverse}(S[a,b))\) we have:
\[
  W[j]=S[b-1-j],\qquad
  W[c,d)=\operatorname{reverse}(S[b-d,b-c)).
\]
The complementing operation leaves the coordinates unchanged, but it changes the binary flag representing whether the slice is complemented or not.

\begin{figure}[H]
\centering
\begin{tikzpicture}[x=1cm,y=1cm,font=\footnotesize]
  \node[anchor=west,font=\bfseries] at (0,4.8)
    {(a) A slice of a slice keeps exact source positions};
  \node[anchor=east] at (2.3,3.9) {$S$};
  \foreach \i/\s in {0/A,1/C,2/G,3/T,4/A,5/A,6/C,7/C,8/G,9/T,10/T,11/A} {
    \node[seq nine cell,fill=black!3] at (2.7+.72*\i,3.9) {\texttt{\s}};
    \node[font=\scriptsize] at (2.7+.72*\i,4.4) {\i};
  }
  \foreach \i/\s in {2/G,3/T,4/A,5/A,6/C,7/C,8/G,9/T}
    \node[seq nine blue] at (2.7+.72*\i,3.9) {\texttt{\s}};
  \foreach \i/\s in {3/T,4/A,5/A,6/C}
    \node[seq nine orange] at (2.7+.72*\i,3.9) {\texttt{\s}};
  \node[anchor=east] at (2.3,2.85) {$V=S[2,10)$};
  \foreach \i/\s in {0/G,1/T,2/A,3/A,4/C,5/C,6/G,7/T} {
    \node[seq nine blue] at (4.14+.72*\i,2.85) {\texttt{\s}};
    \node[font=\scriptsize] at (4.14+.72*\i,2.38) {\i};
  }
  \foreach \i/\s in {1/T,2/A,3/A,4/C}
    \node[seq nine orange] at (4.14+.72*\i,2.85) {\texttt{\s}};
  \draw[seq nine flow,seqNineBlue] (4.14,3.62)--(4.14,3.15);
  \draw[seq nine flow,seqNineBlue] (9.18,3.62)--(9.18,3.15);
  \node[anchor=west,text=seqNineOrange] at (10.35,2.85)
    {$V[1,5)$};

  \node[anchor=west,font=\bfseries] at (0,1.45)
    {(b) Reversal changes which source position comes first};
  \node[anchor=east] at (2.3,.55) {$W$};
  \foreach \i/\s/\p in {0/T/9,1/G/8,2/C/7,3/C/6,4/A/5,5/A/4,6/T/3,7/G/2} {
    \node[seq nine blue] at (4.14+.72*\i,.55) {\texttt{\s}};
    \node[font=\scriptsize] at (4.14+.72*\i,.03) {$j=\i$};
    \node[font=\scriptsize,text=seqNineBlue] at (4.14+.72*\i,1.08) {\p};
  }
  \foreach \i/\s in {1/G,2/C,3/C}
    \node[seq nine orange] at (4.14+.72*\i,.55) {\texttt{\s}};
  \node[anchor=east,font=\scriptsize,text=seqNineBlue] at (3.63,1.08)
    {source position};
  \node[anchor=west,align=left,text=seqNineOrange] at (10.35,.55)
    {$W[1,4)$\\reads $S[6,9)$\\backward: \texttt{GCC}};
\end{tikzpicture}
\caption{Blue cells form the first slice and orange cells show a second slice}
\label{fig:sequence-slice-coordinates}
\end{figure}
